
%=== File containing Global Configuration of the report ===%
%                                                          %
% Copyright (C) ISI - All Rights Reserved                  %
% Proprietary                                              %
% Written by Med Hossam <med.hossam@gmail.com>, April 2016 %
%                                                          %
% @author: HEDHILI Med Houssemeddine                       %
% @linkedin: http://tn.linkedin.com/in/medhossam           %
%==========================================================%

%=========== You MUST type your information here ==========%
% global_config.tex file is designed to configure your     %
% cover pages (main, back and black covers)                %
%==========================================================%

%============= Config new columns type ==============%
\newcolumntype{L}{>{\raggedright\arraybackslash}}
\newcolumntype{R}{>{\raggedleft\arraybackslash}}
\newcolumntype{C}{>{\centering\arraybackslash}}
%==================================================%

%========= Config the cover section ==========%

\title{Design and Implementation of a Debugger with Instruction Trace Capabilities}

\author{Youssef HASNAOUI}
%%% if necessary
% Set isBinomal to true and type second author name
%\setboolean{isBinomal}{true}
%\secondAuthor{Prénom NOM}

\diplomaName{National Diploma in Infotronics Systems Engineering}
\speciality{Microelectronics \& Embedded Systems}
%\speciality{Génie des Télécommunications et Réseaux}
%\speciality{Génie Informatique des Systèmes Industriels}

%% Encadrant professionnel
\proFramerNameOne{Mohamed HAMROUNI}
\proFramerSpecialityOne{Application Engineer}

\academicFramerName{Faten Salem}
\academicFramerSpeciality{Professor}

%% Entreprise d'accueil
\companyName{STMicroelectronics}

%% Année universitaire
\collegeYear{2025 - 2026}

%%%%%% Signatures section %%%%%%

% You can simply remove theses sentences by typing an empty string
% \proSignSentence{}

\proSignSentence{J'autorise l'étudiant à faire le dépôt de son rapport de stage en vue d'une soutenance.}


%%% AR

%% To use latin characters inside the arabic text
% just put them inside the command \textLR{}
%%%%

%%% FR
\frenchAbstract{Ce rapport décrit un débogueur capable de capturer, de décoder et de présenter la trace d'instructions d'un microcontrôleur STM32. La trace d'instructions enregistre le chemin que le programme a réellement suivi dans son code. Le matériel qui la génère équipe une grande partie de la gamme STM32, mais il n'est exploitable qu'au moyen d'outils propriétaires et de sondes dont la plupart des développeurs ne disposent pas. Deux pilotes CoreSight ont été ajoutés à un serveur de débogage libre, afin que l'unité de trace et la mémoire tampon de trace intégrée au circuit puissent être configurées et vidées via la connexion de débogage habituelle. Une chaîne de traitement décode ensuite les octets capturés à l'aide de l'image du programme, reconstruit les instructions exécutées, rattache chacune à une fonction, à un fichier source et à une ligne, et présente le résultat au sein d'un environnement de développement, pendant une session de débogage. Une capture de 128 octets, extraite de la mémoire tampon intégrée au circuit et obtenue sans recourir aux broches de trace, permet de reconstruire 114 instructions exécutées.}
\frenchAbstractKeywords{Trace d'instructions, CoreSight, Embedded Trace Macrocell, Débogage embarqué, STM32, Debug Adapter Protocol, OpenOCD.}
%% dont go over 10 lines
%%% EN
\englishAbstract{This report describes a debugger that captures, decodes and presents the instruction trace of an STM32 microcontroller. Instruction trace records the path a program actually took through its code. The hardware producing it is present across much of the STM32 range, yet it is reachable only through proprietary tools and probes most developers do not have. Two CoreSight drivers were added to an open-source debug server, so that the trace unit and the on-chip trace buffer may be configured and drained over the ordinary debug connection. A pipeline decodes the captured bytes against the program image, reconstructs the executed instructions, attributes each to a function, a source file and a line, and presents the result inside a development tool during a live debug session. A capture of 128 bytes, drained from the on-chip buffer with no trace pins in use, reconstructs 114 executed instructions.}
\englishAbstractKeywords{Instruction trace, CoreSight, Embedded Trace Macrocell, Embedded debugging, STM32, Debug Adapter Protocol, OpenOCD.}
%% if you want to get rid of the company address just set the boolean variable to false
% PS : it's optional
\setboolean{wantToTypeCompanyAddress}{false}

\companyEmail{contact@company.com}
\companyTel{71 111 111}
\companyFax{71 222 222}
\companyAddressAR{نهج بحيرة ملاران - ضفاف البحيرة - تونس}
\companyAddressFR{Rue du Lac Malaren, Les Berges du Lac 1053 Tunis}

%============= Drafting placeholders ==============%
% Visible markers for artwork that is still to be produced.
% Grep for \needfig, \needtab or \needshot to find every outstanding item.
\newcommand{\needartwork}[3]{%
  \begin{center}
  \fbox{\parbox{0.92\linewidth}{%
    \textbf{[#1]}\\[2pt]
    \textit{Caption:} #2\\[2pt]
    \textit{Content:} #3}}
  \end{center}}
\newcommand{\needfig}[2]{\needartwork{FIGURE TO PRODUCE}{#1}{#2}}
\newcommand{\needtab}[2]{\needartwork{TABLE TO PRODUCE}{#1}{#2}}
\newcommand{\needshot}[2]{\needartwork{SCREENSHOT TO CAPTURE}{#1}{#2}}
% Outstanding work on the text itself, as opposed to missing artwork.
\newcommand{\needwork}[1]{%
  \begin{center}
  \fbox{\parbox{0.92\linewidth}{\textbf{[TODO]}\\[2pt] #1}}
  \end{center}}
%==================================================%

%============= Chapter cross-references ==============%
% The class calls \minitoc from inside \@makechapterhead, which leaves
% \@currentlabel holding the number of the preceding section. A \label placed
% after \chapter therefore records that section number instead of the chapter
% number. \chaplabel sets \@currentlabel to the chapter number first.
% isipfe.cls calls \makeatletter and never restores the catcode, and
% tpl/cover_page.tex relies on that, so the previous state is saved and
% restored here rather than forcing \makeatother.
\expandafter\edef\csname restore@at@code\endcsname{%
  \catcode`\noexpand\@=\the\catcode`\@\relax}
\makeatletter
\newcommand{\chaplabel}[1]{%
  \protected@edef\@currentlabel{\thechapter}%
  \label{#1}}
\csname restore@at@code\endcsname
%==================================================%
